\SourceLeaf{040}
\fancyhead[L]{\small DEMONSTRATIO NOVA ALTERA THEOREMATIS}
\begin{center}
\textbf{DEMONSTRATIO NOVA ALTERA THEOREMATIS}\\[12pt]
\textbf{OMNEM FUNCTIONEM ALGEBRAICAM RATIONALEM INTEGRAM}\\[12pt]
UNIUS VARIABILIS\\[12pt]
\textbf{IN FACTORES REALES PRIMI VEL SECUNDI GRADUS RESOLVI POSSE.}\\[12pt]
\rule{3cm}{0.3pt}
\end{center}

\SecondArticle{1} Quamquam demonstratio theorematis de resolutione functionum algebraicarum integrarum in factores, quam in commentatione sedecim abhinc annis promulgata tradidi, tum respectu rigoris tum simplicitatis nihil desiderandum relinquere videatur, tamen haud ingratum fore geometris spero, si iterum ad eandem quaestionem gravissimam revertar, atque e principiis prorsus diversis demonstrationem alteram haud minus rigorosam adstruere coner. Pendet scilicet illa demonstratio prior, partim saltem, a considerationibus geometricis: contra ea, quam hic exponere aggredior, principiis mere analyticis innixa erit. Methodorum analyticarum, per quas usque ad illud quidem tempus alii geometrae theorema nostrum demonstrare susceperunt, insigniores loco citato recensui, et quibus vitiis laborent copiose exposui. Quorum gravissimum ac vere radicale omnibus illis conatibus, perinde ac recentioribus, qui quidem mihi innotuerunt, commune: quod tamen neutiquam inevitabile videri in demonstratione analytica, iam tunc declaravi. Esto iam penes peritos iudicium, an fides olim data per has novas curas plene sit liberata.

\SecondArticle{2} Disquisitioni principali quaedam praeliminares praemittentur, tum ne quid deesse videatur, tum quod ipsa forsan tractatio iis quoque, quae ab aliis iam \SourceLeaf{041}
delibata fuerant, novam qualemcunque lucem affundere poterit. Ac primo quidem de altissimo divisore communi duarum functionum algebraicarum integrarum unius indeterminatae agemus. Ubi praemonendum, hic semper tantum de functionibus integris sermonem esse: e qualibus duabus si productum confletur, utraque huius \textit{divisor} vocatur. Divisoris \textit{ordo} ex exponente summae potestatis indeterminatae quam continet diiudicatur, nulla prorsus coëfficientium numericorum ratione habita. Ceterum quae ad divisores communes functionum pertinent, eo brevius absolvere licet, quod iis, quae ad divisores communes numerorum spectant, omnino sunt analoga.

Propositis duabus functionibus $Y,Y'$ indeterminatae $x$, quarum prior sit ordinis altioris aut saltem non inferioris quam posterior, formabimus aequationes sequentes
\[
\begin{gathered} Y=qY'+Y''\\ Y'=q'Y''+Y'''\\ Y''=q''Y'''+Y''''\\ \text{etc. usque ad}\\ Y^{(\mu-1)}=q^{(\mu-1)}Y^{(\mu)}
\end{gathered}
\] ea scilicet lege, ut primo $Y$ dividatur sueto more per $Y'$; dein $Y'$ per residuum primae divisionis $Y''$, quod erit ordinis inferioris quam $Y'$; tunc rursus residuum primum per secundum $Y'''$ et sic porro, donec ad divisionem absque residuo perveniatur, quod tandem necessario evenire debere inde patet, quod ordo functionum $Y',Y'',Y'''$ etc. continuo decrescit. Quas functiones perinde atque quotientes $q,q',q''$ etc. esse functiones \textit{integras} ipsius $x$, vix opus est monere. His praemissis, manifestum est,

I. regrediendo ab ultima istarum aequationum ad primam, functionem $Y^{(\mu)}$ esse divisorem singularum praecedentium, adeoque certo divisorem communem propositarum $Y,Y'$.

II. Progrediendo a prima aequatione ad ultimam, elucet, quemlibet divisorem communem functionum $Y,Y'$ etiam metiri singulas sequentes, et proin etiam ultimam $Y^{(\mu)}$. Quamobrem functiones $Y,Y'$ habere nequeunt ullum divisorem communem altioris ordinis quam $Y^{(\mu)}$, omnisque divisor communis eiusdem ordinis ut $Y^{(\mu)}$ erit ad hunc in ratione numeri ad numerum, unde hic ipse pro divisore communi summo erit habendus.
